% =====================================================================
%  SECTION 1 — Introduction
%  Draft 1. Written LAST and written against sec11_limitations.tex:
%  every claim here is one that section 11 permits.
% =====================================================================

\section{Introduction}
\label{sec:intro}

Effective properties of heterogeneous materials are computed constantly and
certified rarely. A designer screening candidate microstructures, a
qualification engineer justifying an allowable, and an uncertainty study
propagating microstructural variability all need a number for the effective
conductivity or stiffness of a given arrangement of phases. Established
computational homogenization---finite elements, or the FFT schemes of
Moulinec and Suquet~\cite{Moulinec1998}, supplies that number quickly and
accurately. What it does not supply is an error bar: the output is an
estimate whose distance from the exact effective tensor is unknown without
a separate convergence study, and in our own measurements an FFT solve
returned an estimate $3.867\%$ from the analytic value with nothing in its
output to reveal the discrepancy.

Variational bounds resolve this. The primal and complementary energy
principles bracket the effective tensor between two computable matrices,
and the width of the bracket is itself a certificate of accuracy. The
construction dates to the hypercircle method of Prager and
Synge~\cite{PragerSynge1947} and has been developed into guaranteed
computable bounds for homogenization by several
groups~\cite{Vondrejc2015,Vondrejc2016,Gaynutdinova2023}.
Its neural realisation is recent: Gaynutdinova et
al.~\cite{Gaynutdinova2026} train primal and dual networks for the scalar
conductivity cell problem and certify the resulting bounds by projecting
the networks onto conforming piecewise-linear finite elements, where the
energy integrals can be evaluated exactly. Their formulation is, however,
per-instance: every microstructure requires its own training run, which is
precisely the wrong cost model for the many-query settings above.

Machine-learned surrogates address the cost. Amortized homogenization is by
now well established: graph networks that take the microstructure as an
input and generalise to geometries absent from training~\cite{Hendriks2024},
physics-informed solvers of the cell problem with architecturally enforced
periodicity~\cite{Jiang2023}, and label-free energy-minimising neural
solvers for periodic microstructures~\cite{Xing2025}. Most such approaches
report estimates rather than a geometry-specific computable error bound.
A notable exception is the Voigt--Reuss network of Keshav, Herb and
Fritzen~\cite{Keshav2025}, which predicts stiffness tensors
for unseen microstructures and enforces, by construction, that every
prediction lies inside the analytic Voigt--Reuss envelope in the L\"owner
sense. That envelope, however, is a function of the phase fractions alone:
it assigns the same admissible set to every composite of the same
composition, and therefore certifies physical admissibility rather than
per-instance accuracy.

\medskip
The three positions may be summarised as follows.

\begin{center}
\resizebox{\textwidth}{!}{%
\begin{tabular}{lll}
\toprule
 & certificate & scope \\
\midrule
Gaynutdinova et al.~\cite{Gaynutdinova2026} & variational, tight & per-instance, scalar \\
Voigt--Reuss net~\cite{Keshav2025} & amortized, fraction-only & tensor elasticity \\
\textbf{this work} & \textbf{amortized and per-instance tight} & scalar and tensor, 2D \\
\bottomrule
\end{tabular}}
\end{center}
\medskip

This paper occupies the missing entry. We stress the precise claim:
amortization over unseen microstructures is not itself new~\cite{Hendriks2024,Jiang2023,Xing2025},
and exact-integration certification of neural trial fields is not
new~\cite{Gaynutdinova2026}. Our contribution is their combination,
\emph{certified} amortization: two-sided bounds for microstructures not
used to train the operator, computed at inference cost, under the
certification contract stated below. The results demonstrate tighter
intervals than the analytic envelope on the reported test cases; they do
not establish superiority to every certified or learned competitor.

\subsection{Contributions}

\paragraph{Certified amortization.} We condition primal and dual trial
fields on the microstructure in a way that preserves conformity, so that a
single trained operator emits certified bounds for unseen microstructures.
The reported inference timings, $19$\,ms for conductivity and $62$\,ms for
elasticity, are measured on the controlled generated families used here.
Unlike amortized
surrogates that regress the effective tensor
directly~\cite{Hendriks2024,Keshav2025}, the operator here emits
\emph{admissible fields}, from which the bounds follow by the variational
principles of Section~\ref{sec:primaldual}; the guarantee is therefore a
property of the construction rather than of the training data. Under the
certification contract, the bound theorems hold for every input whether or
not it resembles the training distribution. The interval width is a
geometry-specific measure of trial-field suboptimality, not a universal
detector of distribution shift. We show that this conditioning must be
smooth and periodic
within the automatic-differentiation graph---the natural implementation
using per-point coefficient lookups destroys conformity and voids the
guarantee---a constraint that cannot arise in per-instance formulations.

\paragraph{Certificates that see geometry.} On the controlled scalar
conductivity test cells, certified intervals are between two and four times
narrower than the Voigt--Reuss envelope. More directly, for two
microstructures of identical phase fraction $0.500$, a rank-one laminate and
a checkerboard, the envelope is by construction the same, while the
certified widths are $18.3\%$ and $43.2\%$ respectively. A guarantee
derived from composition cannot make this distinction; a variational
certificate does, at inference cost.

\paragraph{Tensor-valued elasticity.} We extend the amortized construction
to plane-strain elasticity, representing statically admissible periodic
stresses through Airy potentials and certifying them by projection onto
$C^1$ Bogner--Fox--Schmit elements, where global $C^1$ continuity is
exactly what makes the interpolated traction continuous across element
edges. The classical form of this dual is due to Gaynutdinova et
al.~\cite{Gaynutdinova2023}; the contribution here is carrying it onto
amortized neural trial fields with exact-integration certification.

\paragraph{A criterion for collocation adequacy.} That a
collocation-evaluated loss can be arbitrarily misleading is
known~\cite{Berrone2022}, and integration overfitting has been observed
empirically in this
setting~\cite{Gaynutdinova2026}. We contribute prediction and prescription:
a computable comparison between the trained field's spectral bandwidth and
the Nyquist limit of the quadrature rule that flags the failure before
certification degrades, functions as a per-run and label-free accept/reject
gate, and motivates how to choose the collocation density. In the reported
experiments, bandwidth exceeding the Nyquist limit accompanied inaccurate
collocation energies, including reported upper bounds below the analytic
reference. We treat the diagnostic as empirical; it is not a universal
acceptance criterion for physics-informed networks.

\paragraph{Applications.} Because the certificate is differentiable, the
design objective can be the guaranteed value itself. In illustrative
fixed-budget synthetic calculations, optimising inclusion positions raises a
guaranteed conductivity floor by $55.3\%$ and a guaranteed directional
stiffness floor by $36.5\%$.
The construction is self-penalising: a looser dual field yields a lower guaranteed floor, so the optimiser cannot profit from degrading the
surrogate, which is the failure mode that makes surrogate-driven design
untrustworthy. Separately, for the same specified geometry, phase law,
boundary conditions, tensor convention and discretisation, a reference value
falling outside the interval establishes an inconsistency that must be
resolved; we report a case in which the certificate identified an
under-converged finite-volume solution, confirmed subsequently by mesh
refinement.

\paragraph{Intended use.} We state at the outset the regime each mode
serves, since the certified widths differ by an order of magnitude between
them. In the synthetic elasticity experiment, amortized certificates had a
median width of $28\%$ and are intended for \emph{screening}: a requirement falling outside the interval
decides a candidate provably, and Section~\ref{sec:amortized:screening}
shows that this settles roughly two thirds of a design pool at $10$\,ms per
candidate. Per-instance certification, reaching widths below $1\%$, is
intended for final design. We do not claim that a $28\%$ interval is a
design-grade property estimate.

\subsection{Scope}

We state the boundaries here rather than only in
Section~\ref{sec:limitations}. The construction is two-dimensional, linear
and periodic, and excludes voids and rigid phases. In the controlled
families studied here, certified widths degrade towards the analytic envelope
for corner-singular and bicontinuous morphologies above roughly $50{:}1$
phase contrast; we map this boundary explicitly rather than reporting only
favourable cases, and identify the limitation as the approximation of
complex corrector fields by smooth networks rather than any failure of the
certification machinery. We claim
no speed advantage over classical homogenization. An FFT solve of these cells is roughly fifty times faster per query, and the amortized route
becomes preferable only beyond approximately $300$ certified queries. Nor
do we claim experimental validation: the certificate concerns a specified
idealised periodic model and is validated here against closed-form effective
tensors and independent finite-volume, finite-element and FFT calculations
for matching model problems. The reported plane-strain performance results
are otherwise controlled numerical demonstrations on author-generated
periodic two-phase cells. We additionally report a fixed external coverage
audit on $200$ labelled cells from the public DaRUS-5730 plane-strain
dataset~\cite{Keshav2026data}. It uses the released native $100\times100$
rasters, an explicit image/tensor coordinate-frame conversion, and the
released phase law; the latter conversion is independently checked by
periodic $Q_1$ calculations on three formerly discrepant cells. We do not,
however, report a head-to-head trained-predictor comparison. Such a
comparison requires the same fixed split and a documented match of geometry,
phase law, boundary condition, tensor convention and discretisation.

\subsection{Outline}

Section~\ref{sec:primaldual} states the primal--dual certificates and fixes
notation. Section~\ref{sec:admissible} describes how admissibility is
enforced by architecture, and Section~\ref{sec:exact} how the certificates
are made exact by projection onto conforming finite elements.
Section~\ref{sec:amortized} presents the amortized construction and its
results, and Section~\ref{sec:tightness} compares the resulting
certificates against the analytic envelope. Section~\ref{sec:quadrature}
develops the collocation-adequacy criterion, and Section~\ref{sec:limits}
maps the boundary of usefulness across contrast and morphology.
Sections~\ref{sec:applications} and~\ref{sec:cost} give applications and an
honest cost comparison, and Section~\ref{sec:limitations} collects the
limitations.
